\section{Problem Definition and Cost Model}
\label{sec:problem}

\subsection{Setting}
We consider a recurring Spark pipeline that is about to be executed $N$ more times over data
with similar characteristics. A workload $w$ is described by three observable features: its
data volume, the number of files its input is split into (fragmentation), and the volume of
data it shuffles. The operator may leave the pipeline in its \emph{static} configuration or
apply one optimization action $a$ once, before the next execution. We study two actions of
different kinds:
\begin{itemize}
  \item \textbf{A1, shuffle-partition tuning} (execution level): change
        \texttt{spark.sql.shuffle.partitions} from its default of 200 to a value chosen for $w$;
  \item \textbf{A2, file compaction} (storage level): rewrite the input files of $w$ into
        fewer, larger files.
\end{itemize}
The static baseline uses 200 shuffle partitions and no compaction. We restrict the action set
to $\mathcal{A} = \{\textit{none}, \mathrm{A1}, \mathrm{A2}\}$; a combined A1+A2 action is
deliberately excluded so that the effect of each action can be attributed cleanly.

\subsection{Total cost and break-even point}
Let $\Crun(a)$ be the cost of one execution of the pipeline after action $a$ has been applied,
and $\Copt(a)$ the one-time cost of applying it. With $\Cstatic = \Crun(\textit{none})$ and
$\Copt(\textit{none}) = 0$, the total cost of choosing $a$ for the next $N$ executions is
\begin{equation}
  TC(a, N) = \Copt(a) + N \cdot \Crun(a).
  \label{eq:tc}
\end{equation}
An action breaks even over $N$ executions when $TC(a,N) \le TC(\textit{none},N)$.
Writing $\Delta(a) = \Cstatic - \Crun(a)$ for the per-execution saving, the optimization
break-even point (OBEP) is the smallest $N$ at which the accumulated saving $N \cdot \Delta(a)$
covers the one-time cost:
\begin{equation}
  \Nstar(a) =
  \begin{cases}
    \left\lceil \Copt(a) / \Delta(a) \right\rceil & \text{if } \Delta(a) > 0,\\
    \infty & \text{otherwise.}
  \end{cases}
  \label{eq:obep}
\end{equation}
An action with no saving never pays off, however often the pipeline runs; an action whose
one-time cost is smaller than a single saving pays off at the first execution
($\Nstar = 1$).

\subsection{Estimating cost}
All costs are \emph{estimated} from published cloud prices; nothing is billed. We price compute
at the AWS EC2 on-demand rate per vCPU-hour, $p$ \cite{aws2026ec2pricing}, and convert each
measured execution into cost with two proxies:
\begin{align}
  \text{P1 (time-based):}\quad & C = t_{\mathrm{wall}} \cdot v \cdot p / 3600, \label{eq:p1}\\
  \text{P2 (resource-based):}\quad & C = t_{\mathrm{exec}} \cdot p / 3600, \label{eq:p2}
\end{align}
where $t_{\mathrm{wall}}$ is the wall-clock time of the execution in seconds, $v$ the number of
vCPUs reserved for it, and $t_{\mathrm{exec}}$ the executor run time summed over all tasks
(task-seconds). P1 corresponds to paying for a reserved machine for as long as the job runs;
P2 charges only for the core time the tasks actually used. The two can disagree, because
extra parallelism shortens wall-clock time while adding scheduling work to every task, so we
report both.

When $\Copt$ and $\Crun$ are expressed with the same proxy, the price $p$ cancels in
(\ref{eq:obep}): $\Nstar$ depends only on measured times. Prices therefore matter for absolute
cost figures and for the sensitivity analysis, but not for the break-even count itself.

\subsection{One-time cost of each action}
\textbf{A1.} The best partition count is not known in advance; in practice it is found by
trying candidate settings. Our primary definition therefore charges A1 with the cost of a
search: one execution under each of the $K$ candidate settings,
\begin{equation}
  \Copt(\mathrm{A1}) = \sum_{k=1}^{K} C\bigl(\text{run at setting } k\bigr),
  \label{eq:copt-a1}
\end{equation}
treating every search execution as pure overhead. This is conservative: the search runs also
produce pipeline output, so the true overhead is smaller. As a lower bound we also report a
\emph{rule-based} variant in which the setting is chosen from workload features without any
search, so that $\Copt(\mathrm{A1}) \approx 0$ and A1 pays off at the first execution whenever
$\Delta(\mathrm{A1}) > 0$. The gap between the two variants measures what prior knowledge of
a good setting is worth. $\Crun(\mathrm{A1})$ is the per-execution cost at the chosen setting.

\textbf{A2.} $\Copt(\mathrm{A2})$ is the measured cost of the compaction job itself, which
reads the fragmented input and writes the compacted copy. $\Crun(\mathrm{A2})$ is the cost of
one pipeline execution over the compacted data.

\subsection{Assumptions}
The model assumes that (i) the workload's characteristics stay stable over the $N$ executions,
so a saving measured once recurs; (ii) $N$ is known or can be estimated from the pipeline's
schedule; (iii) compaction is not undone by new writes during the horizon; and (iv) the cost
of evaluating the decision rule itself is negligible. In our study $N$ is a parameter of the
model and is not physically executed; we evaluate $N \in \{1, 5, 10, 50\}$.
